\begin{pastexam}[label={px:ncku-111-cce-ds-05}]{成大電通所 111 學年度 第 5 題}
（追蹤）

(10 \textbf{pts}) Given an input array $L$=\{61, 127, 15, 7, 11, 10, 5, 51, 3\}, where
\textit{left} and \textit{right} denote the indexes of the first and last elements in the array
(i.e., \textit{left}= 0 and \textit{right} =8).
\begin{enumerate}[label=\alph*.]
\item (5 \textbf{pts}) Quick sort algorithm uses the divide an conquer strategy to sort number. To
avoid imbalanced partition, we find the pivot in the location estimated by the function f() and
swap it with the first element in the sub-array in each iteration, where the f() function is
shown as follows:
\begin{center}f(\textit{right})= 4* \textit{right} mod $N$,\end{center}
where $N$ is the number of elements in a sub-array (i.e., $N$=9 in the beginning)\\
Please show the resulting arrays after the first and the second iterations in the recursive
procedure.
\item (5 \textbf{pts}) Use a 4-basd radix sort to rearrange numbers in the array. Please show the
result in a bucket list and the resulting array in each iteration.
\end{enumerate}
\end{pastexam}

\begin{pastsoln}
答案：(a) 第一次分割後是 7, 3, 5, 10, 11, 61, 15, 51, 127，第二次後是
3, 5, 7, 10, 11, 61, 15, 51, 127；(b) 四趟，見下表。 (a) 屬資結書第~14~章，(b) 屬第~15~章，
本題放在第~14~章。

\textbf{(a) 快速排序。} 題目沒有指定分割的寫法，這裡用基準在最左邊、兩個指標從兩端往中間掃、
最後把基準換到中間的寫法（Horowitz 等人的資料結構課本就是這種），每一次「iteration」指一次分割；第二次是
遞迴先處理的左半段。 $f$ 算出來的位置照題意直接當陣列索引；前兩次分割的子陣列都從索引 $0$
開始，看成絕對位置或相對位置結果一樣。

\begin{pseudo}{本題採用的分割}
\Function{Partition-First}{$A, \mathit{left}, \mathit{right}$}
    \State $N \gets \mathit{right} - \mathit{left} + 1$；$p \gets (4 \cdot \mathit{right}) \bmod N$
    \State 交換 $A[\mathit{left}]$ 和 $A[p]$；$x \gets A[\mathit{left}]$ \Comment{基準}
    \State $i \gets \mathit{left}$；$j \gets \mathit{right} + 1$
    \While{\TRUE}
        \State \textbf{repeat} $i \gets i + 1$ \textbf{until} $i > \mathit{right}$ \OR $A[i] \ge x$
        \State \textbf{repeat} $j \gets j - 1$ \textbf{until} $A[j] \le x$
        \If{$i \ge j$}
            \State \textbf{break}
        \EndIf
        \State 交換 $A[i]$ 和 $A[j]$
    \EndWhile
    \State 交換 $A[\mathit{left}]$ 和 $A[j]$；\Return $j$
\EndFunction
\end{pseudo}

\textbf{第一次}（$\mathit{left} = 0$，$\mathit{right} = 8$，$N = 9$）：$f = 4 \cdot 8 \bmod 9 =
32 \bmod 9 = 5$，基準是 $L[5] = 10$，和 $L[0]$ 交換。
\begin{center}
\small
\begin{tabular}{l*{9}{c}}
\toprule
索引 & 0 & 1 & 2 & 3 & 4 & 5 & 6 & 7 & 8 \\
\midrule
開始 & 61 & 127 & 15 & 7 & 11 & 10 & 5 & 51 & 3 \\
基準換到最前面 & 10 & 127 & 15 & 7 & 11 & 61 & 5 & 51 & 3 \\
$i = 1, j = 8$，交換 & 10 & \textbf{3} & 15 & 7 & 11 & 61 & 5 & 51 & \textbf{127} \\
$i = 2, j = 6$，交換 & 10 & 3 & \textbf{5} & 7 & 11 & 61 & \textbf{15} & 51 & 127 \\
$i = 4, j = 3$，$i \ge j$ 停 & 10 & 3 & 5 & 7 & 11 & 61 & 15 & 51 & 127 \\
基準和 $A[3]$ 交換 & \textbf{7} & 3 & 5 & \textbf{10} & 11 & 61 & 15 & 51 & 127 \\
\bottomrule
\end{tabular}
\end{center}
第一次之後：$[7, 3, 5, 10, 11, 61, 15, 51, 127]$，基準 $10$ 在索引 $3$，左邊都比它小、右邊都比它大。

\textbf{第二次}，遞迴處理左半段 $A[0..2] = [7, 3, 5]$：$N = 3$，$f = 4 \cdot 2 \bmod 3 = 8 \bmod 3 =
2$，基準是 $A[2] = 5$，換到最前面得 $[5, 3, 7]$。 $i$ 跳過 $3$ 停在索引 $2$（$7 \ge 5$），$j$ 從
索引 $2$ 往左停在索引 $1$（$3 \le 5$），$i \ge j$ 停；基準和 $A[1]$ 交換，得 $[3, 5, 7]$。 整個
陣列：
\[
[3, 5, 7, 10, 11, 61, 15, 51, 127].
\]

\textbf{(b) 以 $4$ 為基數的 LSD 基數排序}（資結書第~15.2.2~節）。 最大的 $127 = 1333_4$，所以每個數
寫成 $4$ 位的四進位，做 $4$ 趟，每趟依一位數放進 $0$--$3$ 號桶子（保持先來後到），再依桶子順序
收回：
\begin{center}
\small
\begin{tabular}{*{9}{c}}
\toprule
61 & 127 & 15 & 7 & 11 & 10 & 5 & 51 & 3 \\
\midrule
$0331$ & $1333$ & $0033$ & $0013$ & $0023$ & $0022$ & $0011$ & $0303$ & $0003$ \\
\bottomrule
\end{tabular}
\end{center}
\begin{center}
\small
\begin{tabular}{lllll}
\toprule
 & 第 1 趟（$4^0$ 位） & 第 2 趟（$4^1$ 位） & 第 3 趟（$4^2$ 位） & 第 4 趟（$4^3$ 位） \\
\midrule
桶 0 & & 51, 3 & 3, 5, 7, 10, 11, 15 & 3, 5, 7, 10, 11, 15, 51, 61 \\
桶 1 & 61, 5 & 5, 7 & & 127 \\
桶 2 & 10 & 10, 11 & & \\
桶 3 & 127, 15, 7, 11, 51, 3 & 61, 127, 15 & 51, 61, 127 & \\
\bottomrule
\end{tabular}
\end{center}
每一趟收回後的陣列：
\begin{itemize}
\item 第 1 趟：$61, 5, 10, 127, 15, 7, 11, 51, 3$
\item 第 2 趟：$51, 3, 5, 7, 10, 11, 61, 127, 15$
\item 第 3 趟：$3, 5, 7, 10, 11, 15, 51, 61, 127$
\item 第 4 趟：$3, 5, 7, 10, 11, 15, 51, 61, 127$
\end{itemize}
第 3 趟後其實已經排好（只有 $127 \ge 4^3$），但 $127$ 有 $4$ 位，第 4 趟還是要做才保證正確。 時間
$\Theta(d(n + r)) = \Theta(4 \cdot (9 + 4))$。
\end{pastsoln}
